\section{Difficulty Adjustment via Amplitude Amplification}


\subsection{The Difficulty Search Problem}

In the consensus protocol, each proposer must find a truncation
parameter $N^*$ and evaluation point $z^*$ such that the computed
Eisenstein value $\tilde{G}_k^{(N^*)}(z^*)$ satisfies a difficulty
target $\tau$:

\[
\left| \tilde{G}_k^{(N^*)}(z^*) - \tau \right| < \Delta
\]

where $\Delta$ is the difficulty window. Classically, searching over
$M$ candidate pairs $(N, z)$ requires $O(M)$ evaluations.

\subsection{Grover-Accelerated Search}

Amplitude amplification (Grover's algorithm~\cite{grover1996})
provides a quadratic speedup. Define the oracle
$O_\tau: \ket{N,z}\ket{b} \rightarrow \ket{N,z}\ket{b \oplus f(N,z)}$
where $f(N,z) = 1$ iff $|\tilde{G}_k^{(N)}(z) - \tau| < \Delta$.

\begin{proposition}[Difficulty Search Complexity]
Given $M$ candidate parameter pairs with $t$ satisfying solutions,
the quantum difficulty search finds a valid $(N^*, z^*)$ in
$O(\sqrt{M/t})$ oracle queries, each requiring one Eisenstein
evaluation circuit.
\end{proposition}

The protocol dynamically adjusts the difficulty window $\Delta$ based
on block production rate, analogous to Bitcoin's difficulty retargeting
but with the search itself benefiting from quantum speedup.

\subsection{Classical Verification}

A critical property: while \emph{finding} a valid $(N^*, z^*)$ benefits
from quantum speedup, \emph{verifying} a proposed solution requires only
a single classical Eisenstein evaluation. This asymmetry ensures that
quantum-equipped proposers gain an advantage in block production without
preventing classical nodes from validating the chain.

%% ============================================================
